\section*{Chapter \ref{ch:s1:index-rebalancing} --- Index Rebalancing}

\begin{solution}{exo:s1:index-rebalancing:1}
A non-member joins at rank 285 or better ($300 - 15$); a member leaves beyond rank 315 ($300 + 15$).
\end{solution}

\begin{solution}{exo:s1:index-rebalancing:2}
The funds must buy $0.15/0.006 = 25$ days of volume: $0.7 \times 0.02 \times \sqrt{25} = 7.0\%$.
\end{solution}

\begin{solution}{exo:s1:index-rebalancing:3}
The announcement trade earns the half not taken early, $0.5 \times 7.0\% = 3.5\%$; the liquidity provider earns the half that reverses, also 3.5\%, before
costs and before the stock's own movements.
\end{solution}

\begin{solution}{exo:s1:index-rebalancing:4}
A stock near the cut-off must cross the band to change status, so small moves around the threshold do not flip it in and out: fewer changes, less trading
by the \omterm{def:s1:index-rebalancing:fund}{index funds}. Changes become rarer and involve stocks that moved further, which makes them easier to predict at the \omterm{def:s1:index-rebalancing:fund}{rank day} but a band also makes the
outcome depend on the stock's current membership.
\end{solution}

\begin{solution}{exo:s1:index-rebalancing:5}
The \omterm{def:s1:index-rebalancing:fund}{index funds}' demand, and so the total push, is the same; arbitrage capital only changes when the price moves. The more capital predicts the change and
buys early, the more of the push happens before the announcement and the less is left between the announcement and the effective close.
\end{solution}

\begin{solution}{exo:s1:index-rebalancing:6}
About a quarter of its positions are in stocks that do not join and earn no push (and may be hurt by the unwinding of others' positions); it also misses
more than half of the actual additions. Its return per position is well below the perfect predictor's, and its risk higher.
\end{solution}

\begin{solution}{exo:s1:index-rebalancing:7}
With 30\% of shares tracked and only a fifth taken early, the announcement trade earns 5.7\% per event; with 5\% tracked and four fifths taken early, 0.4\%.
The first has a large push (7.5\% on average) with most of it still to come; the second a small push (3.1\%) mostly gone before the announcement.
\end{solution}

\begin{solution}{exo:s1:index-rebalancing:8}
The event has changed: more capital now predicts the changes, so more of the move happens before the announcement, and the part after it has shrunk;
Patel and Welch find reversion in the 2000s. Measure the event on recent years, with costs and the stocks' own drifts, and ask who else is buying.
\end{solution}

\begin{solution}{pb:s1:index-rebalancing:1}
\begin{enumerate}
\item A fund holding an index in its weights; the date on which a rules-based index measures membership criteria; the probability, before the \omterm{def:s1:index-rebalancing:fund}{rank day}, that a
stock will be a member after the reconstitution.
\item Semi-annual from 2026: June (\omterm{def:s1:index-rebalancing:fund}{rank day} April 30, effective after the close on June 26) and December (\omterm{def:s1:index-rebalancing:fund}{rank day} October 30, effective after the close on
December 11).
\item The 300 largest stocks, reconstituted every 126 days with a band of 15 ranks; announcement 16 days and effective close 40 days after the \omterm{def:s1:index-rebalancing:fund}{rank day};
331 additions and 270 deletions over seventeen reconstitutions.
\item It keeps stocks near the cut-off from churning in and out.
\item By simulating each stock's capitalisation forward with its own volatility and applying the rule to each draw.
\item 100\% on the \omterm{def:s1:index-rebalancing:fund}{rank day}; 86\% and 72\% for additions 20 days before (deletions 86\% and 74\%); 76\% and 45\% sixty days before (76\% and 46\%).
\item Capitalisations have more time to move.
\item Early predictions leave more of the move ahead but are more often wrong.
\item A share of the push before the announcement, the rest to the effective close, half of it reversing over the next 20 days.
\item 3.75\%, 2.40\% and 2.65\% per event at 15\% tracked and an early share of one half.
\item A significant abnormal return at the announcement of S\&P 500 additions, related to index-fund buying; a permanent rise for additions but no permanent
fall for deletions.
\item Stocks that grow into an index are recent winners, and the synthetic market plants momentum.
\item \textbf{Named result.} Membership changes are predicted with a precision of 86\% and a recall of 72\% twenty days before the \omterm{def:s1:index-rebalancing:fund}{rank day} (76\% and 45\%
sixty days before); the announcement trade earns 2.2\% to 5.7\% per event as the tracked share rises from 5\% to 30\% with a fifth taken early, and 0.4\% to
1.3\% with four fifths taken early.
\item Reversion in the 2000s: no permanent demand shift for S\&P 500 changes.
\item Earlier, to predictors, and later, to liquidity providers at the effective close.
\item Small capital: effective-close liquidity provision in a few names; large: predictions across the whole list.
\item The closing-auction imbalances, the \omterm{def:s1:index-rebalancing:fund}{index funds}' expected demand, and the prices of the most crowded additions.
\item Two smaller events a year instead of one: less demand per event, more events, and a busier calendar for predictors.
\item Its date, size and direction are known in advance.
\item A known, temporary push, most of which is traded before and after the date by those who predict it and those who meet it.
\end{enumerate}
\end{solution}

\begin{solution}{iq:s1:index-rebalancing:1}
\omterm{def:s1:index-rebalancing:fund}{Index funds} must buy it at once, and the other side of that demand is not perfectly elastic: the price rises until sellers are paid enough to supply the
shares. Part of the rise can be information or awareness; the rest is pressure that tends to revert.
\end{solution}

\begin{solution}{iq:s1:index-rebalancing:2}
Predict the changes from the rules and capitalisations before the \omterm{def:s1:index-rebalancing:fund}{rank day}, trade the likely additions and deletions ahead of the preliminary list, and
provide liquidity at the effective close against the \omterm{def:s1:index-rebalancing:fund}{index funds}' demand; size by the expected demand relative to volume and by the confidence of the
prediction.
\end{solution}

\begin{solution}{iq:s1:index-rebalancing:3}
Capitalisation and float as of today, each stock's volatility, the index rules including bands and screens, and the calendar; simulate capitalisations to the
\omterm{def:s1:index-rebalancing:fund}{rank day} to get a probability per stock. The uncertainty grows with the time to the \omterm{def:s1:index-rebalancing:fund}{rank day} and is largest for stocks near the cut-off.
\end{solution}

\begin{solution}{iq:s1:index-rebalancing:4}
Wrong predictions; a crowded trade that has already priced the change; \omterm{def:s1:index-rebalancing:fund}{index funds} that trade before or after the close; a market fall before the \omterm{def:s1:index-rebalancing:fund}{rank day}
that moves the cut-off; stocks near the cut-off that are small and illiquid.
\end{solution}

\begin{solution}{iq:s1:index-rebalancing:5}
If the imbalance is mechanical index demand and the price already reflects heavy pre-positioning, sell into it and plan to buy back as the pressure reverts;
size to the auction's depth and to how much of the move has happened before the close.
\end{solution}

\begin{solution}{iq:s1:index-rebalancing:6}
If arbitrageurs buy $A$ before the announcement and \omterm{def:s1:index-rebalancing:fund}{index funds} buy $D$ at the effective close, the market absorbs $A$ early and $D - A$ later, and sells $A$
back to the funds; with pushes proportional to net demand the early move is proportional to $A$ and the post-announcement move to $D - A$, so the return
after the announcement is proportional to $1 - A/D$ and vanishes as arbitrage capital approaches the index demand.
\end{solution}
